\section{Methodology}
This section introduces \textit{DGLight}, a two-stage reinforcement-learning pipeline for adapting a pretrained large language model to traffic signal control. In the first stage, a CoLight-based DQN critic is trained to estimate action values for signal-phase decisions when provided with structured environment states. In the second stage, the frozen critic scores candidate LLM responses during rollout and the scores are used as the reward signal for GRPO-based policy optimization. Refer to Figure~\ref{fig:dglight_pipeline} for the pipeline architecture.

\begin{figure}[t]
    \centering
    \includegraphics[width=1\linewidth]{dglight-archi.png}
    \caption{Overview of the two-stage \textit{DGLight} pipeline.}
    \label{fig:dglight_pipeline}
\end{figure}

\subsection{Pretraining the CoLight-Based DQN Critic}
The critic used in \textit{DGLight} is based on \textit{CoLight} \cite{10.1145/3357384.3357902}, a graph-attentive deep Q-learning model for network-level traffic signal control. \textit{CoLight} is employed as a pretrained value estimator so that the subsequent GRPO stage can query a traffic-aware action-value function rather than rely exclusively on the raw environment reward. Given the current environment state, the critic predicts action values for each intersection and thereby provides a dense estimate of the quality of candidate signal phases under the prevailing traffic situation.

The \textit{CoLight} critic is composed of three main components. The first component is an observation embedding module that transforms the local traffic state of each intersection into a latent feature representation. The second component consists of one or more multi-head graph-attention layers, which propagate information across neighboring intersections and produce attention-weighted network representations. The final component is a Q-value prediction layer that maps the aggregated features to one scalar value for each admissible signal phase. This architecture combines local state encoding, inter-intersection coordination, and value-based action estimation within a single model.

% \begin{figure}[t]
%     \centering
%     \fbox{
%         \parbox{0.88\linewidth}{
%             \textbf{Critic architecture placeholder.}

%             Replace this box with a figure showing:
%             local traffic observations $\rightarrow$ observation embedding layers
%             $\rightarrow$ graph attention over neighboring intersections
%             $\rightarrow$ per-phase Q-value prediction head.
%         }
%     }
%     \caption{Mockup placeholder for the \textit{CoLight}-based DQN critic architecture used in \textit{DGLight}.}
%     \label{fig:dglight_critic_architecture}
% \end{figure}

Training follows the standard value-based reinforcement learning procedure. The model interacts with the traffic environment to collect state transitions, actions, rewards, and subsequent states ($(s,a,r,s')$), where the stepwise reward is defined to be inversely proportional to queue lengths at each intersection. These transitions are stored in replay memory for off-policy updates. The optimization objective minimizes the Bellman error
\[
\mathcal{L}_{\mathrm{Bellman}}(\theta)
=
\mathbb{E}_{(s,a,r,s')\sim \mathcal{D}}
\left[
\left(
r + \gamma \max_{a'} Q_{\theta^-}(s',a') - Q_{\theta}(s,a)
\right)^2
\right],
\]
where $\mathcal{D}$ denotes the replay buffer, $\theta$ denotes the online network parameters, and $\theta^-$ denotes the target network parameters. Experience replay improves sample efficiency, and the target network improves optimization stability.

\textbf{Disclaimer:} The critic model, its training objective, and its overall learning logic are not original contributions of this dissertation; they are inherited from prior \textit{CoLight} research. My contribution to this component is limited to re-implementing the model in PyTorch for this codebase, whereas the original \textit{CoLight} implementation was based on TensorFlow.

\subsection{LLM Prompting and Action Sampling}
At each signal-switching step $t$, each intersection $i$ is converted into a textual prompt $x_{i,t}$ that summarizes the current traffic condition together with the task instruction and the available signal phases. The prompt is constructed from the environment state used by the controller including the current phase and queue information. Task description is formatted so that the model produces both human-readable reasoning and a final phase choice. A sample prompt is provided in Appendix~\ref{appendix:sample_prompt}.

Given prompt $x_{i,t}$, the policy $\pi_\theta$ samples a group of $k$ candidate responses
\[
\{y_{i,t}^{(1)}, y_{i,t}^{(2)}, \dots, y_{i,t}^{(k)}\}.
\]
Each response is then mapped to a signal phase through a deterministic parser:
\[
a_{i,t}^{(j)} = \operatorname{Parse}(y_{i,t}^{(j)}), \qquad j = 1, \dots, k.
\]
If the response fails to specify a valid action in the predefined action space, it is treated as an invalid output and assigned a fallback reward during training. This allows the method to preserve free-form reasoning generation while still enforcing a discrete control interface at execution time.

\subsection{DQN-Guided Action Scoring in Rollout Collection}
The central idea of \textit{DGLight} is to replace discounted raw environment rewards with action value estimates from the pretrained DQN critic.

\textbf{Stepwise action value assignment.} For the current traffic state at intersection $i$, the frozen critic produces a vector of estimated returns over all admissible signal phases:
\[
\mathbf{q}_{i,t} = \bigl[ Q_\phi(s_{i,t}, a) \bigr]_{a \in A},
\]
where $s_{i,t}$ is the structured traffic state and $Q_\phi$ is the pretrained DQN critic. The critic therefore acts as a traffic-aware evaluator that maps the current state to a vector of action values that reflect both spatial and temporal traffic context.

Each sampled LLM response is scored by the critic through its parsed action:
\[
r_{i,t}^{(j)} =
\begin{cases}
Q_\phi\!\left(s_{i,t}, a_{i,t}^{(j)}\right), & \text{if } a_{i,t}^{(j)} \in A, \\
r_{\mathrm{invalid}}, & \text{otherwise},
\end{cases}
\]
where $r_{\mathrm{invalid}}$ is a fixed fallback reward for malformed or out-of-range decisions. 

\textbf{Environment stepping.} Among the sampled candidates, the action actually executed in the environment is the highest-scoring sampled action,
\[
j_{i,t}^{*} = \arg\max_{j \in \{1,\dots,k\}} r_{i,t}^{(j)},
\qquad
\hat a_{i,t} = a_{i,t}^{(j_{i,t}^{*})}.
\]
Thus, the environment rollout is still driven by sampled LLM outputs, but execution is guided by the critic's action-value estimate.

This design has a major advantage that the critic-derived reward incorporates richer spatial and temporal information than short-horizon cumulative rewards computed directly from the environment. Because the CoLight-based critic aggregates information from neighboring intersections through graph attention and is trained with Bellman backups under a discount factor, $Q_\phi(s_{i,t}, a)$ reflects both network interactions and the expected downstream consequences of choosing action $a$. 

\subsection{GRPO-Based Policy Optimization}
After rollout collection, each intersection-step pair contributes a training record consisting of the prompt and the full critic score vector,
\[
\mathcal{D}_{i,t} = \bigl(x_{i,t}, \mathbf{q}_{i,t}\bigr).
\]
Importantly, the stored supervision is not just the score of the executed action. Instead, the entire per-action vector is retained so that any newly sampled response for the same prompt can be rewarded according to the action it proposes. During GRPO training, a completion $y$ sampled for prompt $x_{i,t}$ is assigned reward
\[
R\!\left(x_{i,t}, y; \mathbf{q}_{i,t}\right) =
\begin{cases}
\mathbf{q}_{i,t}[a(y)], & \text{if } a(y) \in A, \\
r_{\mathrm{invalid}}, & \text{otherwise},
\end{cases}
\]
where $a(y)$ denotes the parsed action of completion $y$. The resulting reward values are then normalized within each sampled group, producing relative advantages that drive the GRPO update.

This optimizer is particularly well suited to the present setting because it provides a natural way to train an LLM using reward signals from an already established value-based traffic-control critic. For a single traffic prompt, multiple completions can be sampled efficiently, parsed into candidate signal phases, and scored by the frozen DQN critic, after which GRPO uses the relative critic rewards within the sampled group to improve the policy. This way, critic-guided GRPO transfers decision-making knowledge from the value-based controller into the LLM. Although the critic is trained on traffic data from a specific city, the LLM is not fine-tuned to overfit any single city. Rather, it is trained to have stronger reasoning patterns for mapping textual traffic states to better signal actions. This strengthens the model's traffic decision-making capacity while preserving transferability.

\subsection{End-to-End DGLight Training Procedure}
Algorithm~\ref{alg:dglight_training} summarizes the complete \textit{DGLight} training loop.

\begin{algorithm}[h]
    \caption{End-to-end \textit{DGLight} training}
    \label{alg:dglight_training}
    \begin{algorithmic}[1]
        \Require traffic environment $\mathcal{E}$, pretrained LLM policy $\pi_\theta$, critic architecture $Q_\phi$, candidate count $k$, invalid-action reward $r_{\mathrm{invalid}}$
        \State Train the CoLight-based DQN critic $Q_\phi$ on traffic transitions
        \State Freeze critic parameters $\phi$
        \For{each policy-training episode}
            \State Reset environment $\mathcal{E}$
            \State Initialize empty GRPO record set $\mathcal{D}_{\mathrm{GRPO}}$
            \While{the rollout episode is not finished}
                \State Initialize empty executed action set $\hat{\mathbf{a}}_t$
                \For{each active intersection $i$}
                    \State Construct textual prompt $x_{i,t}$ from the current traffic state $s_{i,t}$
                    \State Sample candidate responses $\{y_{i,t}^{(1)}, \dots, y_{i,t}^{(k)}\} \sim \pi_\theta(\cdot \mid x_{i,t})$
                    \State Parse each response into a signal phase $a_{i,t}^{(j)} \gets \operatorname{Parse}(y_{i,t}^{(j)})$
                    \State Query the frozen critic to obtain per-action values $\mathbf{q}_{i,t} \gets [Q_\phi(s_{i,t}, a)]_{a \in A}$
                    \For{each sampled response $j \in \{1,\dots,k\}$}
                        \If{$a_{i,t}^{(j)} \in A$}
                            \State $r_{i,t}^{(j)} \gets Q_\phi(s_{i,t}, a_{i,t}^{(j)})$
                        \Else
                            \State $r_{i,t}^{(j)} \gets r_{\mathrm{invalid}}$
                        \EndIf
                    \EndFor
                    \State $j_{i,t}^{*} \gets \arg\max_j r_{i,t}^{(j)}$
                    \State $\hat{a}_{i,t} \gets a_{i,t}^{(j_{i,t}^{*})}$ and add $\hat{a}_{i,t}$ to $\hat{\mathbf{a}}_t$
                    \State Store $(x_{i,t}, \mathbf{q}_{i,t})$ in $\mathcal{D}_{\mathrm{GRPO}}$
                \EndFor
                \State Step the environment with joint action set $\hat{\mathbf{a}}_t$
            \EndWhile
            \State Fine-tune policy $\pi_\theta$ on $\mathcal{D}_{\mathrm{GRPO}}$ with GRPO
        \EndFor
        \State \Return adapted policy $\pi_\theta$
    \end{algorithmic}
\end{algorithm}

In practice, the rollout collection and GRPO optimization stages are decoupled. A full rollout episode is first collected with the frozen critic, and the resulting prompt-score records are then used as a separate training set for the subsequent GRPO update. This separation enables finer-grained optimization, since multiple manual updates can be performed on the same episodic rollout data.

\subsection{Comparison Variant}
For completeness, I also implemented a \textit{Joint-Scored GRPO (JS-GRPO)} variant as a comparison method. Instead of using direct critic $Q$-values to score sampled actions, that variant estimates action quality indirectly from short-horizon joint rollout returns over snapshot-forked simulations, and then projects those returns back to per-intersection action scores for GRPO training. Because this variant is included only to test whether indirect joint scoring can outperform direct critic guidance, I do not treat it as a co-equal methodological contribution in the main text. Its detailed design is provided separately in Appendix~\ref{appendix:joint_scored_grpo}.
